\documentclass[11pt,letterpaper]{article}
\usepackage[T1]{fontenc}
\usepackage{lmodern}
\usepackage[margin=0.83in,headheight=15pt]{geometry}
\usepackage{amsmath,amssymb,amsthm,mathtools}
\usepackage{booktabs,array,longtable,microtype}
\usepackage{xcolor,enumitem,fancyhdr}
\usepackage[colorlinks=true,linkcolor=teal!65!black,citecolor=teal!65!black,urlcolor=blue!55!black,breaklinks=true]{hyperref}
\definecolor{ink}{RGB}{27,43,62}
\pagestyle{fancy}
\fancyhf{}
\fancyhead[L]{\small Shared arithmetic certificates}
\fancyhead[R]{\small ProveIt research}
\fancyfoot[C]{\thepage}
\setlength{\parskip}{4pt}
\setlength{\parindent}{0pt}
\setlength{\emergencystretch}{3em}
\setlist{topsep=4pt,itemsep=2pt}
\setcounter{tocdepth}{2}
\newtheorem{theorem}{Theorem}[section]
\newtheorem{lemma}[theorem]{Lemma}
\newtheorem{proposition}[theorem]{Proposition}
\newtheorem{corollary}[theorem]{Corollary}
\newcommand{\N}{\mathbb N}
\newcommand{\Z}{\mathbb Z}
\newcommand{\Con}{\operatorname{Con}}
\newcommand{\ZFC}{\mathrm{ZFC}}
\newcommand{\file}[1]{\nolinkurl{#1}}
\newcommand{\repoPin}{\texttt{a88aefd7ddab4771534bb049c67350e397c91b9a}}
\newcommand{\NatVars}{667}
\newcommand{\NatRows}{535}
\newcommand{\NatTerms}{23,747}
\newcommand{\NatHeight}{2,649,643,818,040}
\newcommand{\IntVars}{1690}
\newcommand{\IntRows}{842}
\newcommand{\IntTerms}{21,381}
\newcommand{\IntHeight}{2,649,643,805,758}

\hypersetup{pdftitle={Shared arithmetic certificates for Diophantine consistency equations},
 pdfauthor={AI-assisted research for Vladimir Reshetnikov's ProveIt project},
 pdfsubject={Explicit fixed quartics and a 13-integer polynomial for syntactic ZFC consistency}}
\title{\color{ink}\LARGE Shared arithmetic certificates for\\
 Diophantine consistency equations\\[7pt]
 \large A smaller direct quartic and a degree reduction\\
 for the thirteen-integer construction}
\author{Prepared for Vladimir Reshetnikov's ProveIt project\\
 \small AI-assisted mathematical research with exact reproducible artifacts}
\date{23 October 2026\\[3pt]\small Comparison snapshot: \repoPin}
\begin{document}
\maketitle
\begin{abstract}
We construct a fixed quartic polynomial in \NatVars\ nonnegative integer
unknowns, supplied both as \NatRows\ quadratic residuals and as a fully
collected list of \NatTerms\ monomials. Its maximum absolute coefficient is
\NatHeight. Existence of a zero is equivalent to existence of a finite
ZFC refutation. The construction encodes the numerical proof calculus
$C_{222}$ directly, using neither a K215 execution tableau nor Jones's
universal polynomial. Relative to the incoming 3,614-variable direct
quartic, the reduction gives every binomial test one common radix,
all positive-power tests a common Pell certificate, and equal
exponents shared Pell coordinates. Exact nonnegative polynomial
substitutions then eliminate unnecessary coordinates.

The substantive arithmetic step is a simultaneous-completeness theorem.
A single dominating parameter bounds all proposed power outputs; a
finite matrix period supplies one Pell solution divisible by the square
of the product of all required ordinates. Thus independent local
certificates can be replaced by compatible shared witnesses.
An unrestricted-integer quartic is also supplied. A separate
thirteen-integer polynomial has degree $1072$, down from $1088$, at the
same $195$ arithmetic operations. All new equivalences have ordinary
proofs formalizable in ZFC. The archive contains the literal equations,
builders, exact algebra certificates, verification receipts and source
dependencies. No global minimum or new end-to-end proof-assistant
certificate is claimed.
\end{abstract}
\clearpage
{\small\setlength{\parskip}{0pt}\tableofcontents}
\clearpage

\section{The result and its exact interpretation}
\label{sec:result}

Let $\operatorname{Ref}_{\ZFC}$ mean that there exists a finite derivation
of contradiction from the fixed syntactic ZFC axiom predicate. Thus
$\Con(\ZFC)=\neg\operatorname{Ref}_{\ZFC}$, with the usual coding
of finite proofs. Our goal is a fixed explicit polynomial $P$ for which
\[
 \ZFC\vdash
 \left(\neg\exists\mathbf x\in D^n\ P(\mathbf x)=0\right)
 \longleftrightarrow \Con(\ZFC).
\]
The domain $D$ is part of the assertion. Throughout this report
$\N=\{0,1,2,\ldots\}$. Strictly positive variables, natural variables
and unrestricted integers are not interchangeable without a proved
translation.

\begin{theorem}[Delivered equations]
\label{thm:main}
The literal polynomials specified in this archive satisfy
\begin{align*}
 \ZFC&\vdash
 \bigl(\exists\mathbf x\in\N^{\NatVars}\ P_{\mathrm{batch}}(\mathbf x)=0\bigr)
       \longleftrightarrow\operatorname{Ref}_{\ZFC},\\
 \ZFC&\vdash
 \bigl(\exists\mathbf z\in\Z^{\IntVars}\ P_{\mathrm{integer}}(\mathbf z)=0\bigr)
       \longleftrightarrow\operatorname{Ref}_{\ZFC},\\
 \ZFC&\vdash
 \bigl(\exists\mathbf u\in\Z^{13}\ Q_{1072}(\mathbf u)=0\bigr)
       \longleftrightarrow\operatorname{Ref}_{\ZFC}.
\end{align*}
The two direct polynomials have exact degree four. The thirteen-integer
polynomial has exact degree $1072$. Negating these equivalences gives the
requested consistency formulations.
\end{theorem}

This is an ordinary mathematical theorem obtained from the arguments
below and the explicitly identified inherited calculus and arithmetic
theorems. It assumes neither consistency nor arithmetic soundness.
The unbounded mathematical arguments are separate from the finite
software checks. A complete new Lean or object-language ZFC derivation
of the emitted coefficient dictionary has not been produced.

\subsection{What has improved}

The current repository already supplies nine-positive-variable and
ten-integer equations in its high-degree families, as well as a
thirty-seven-integer quartic \cite{latest,wave11}. Our direct construction
addresses a different useful comparison: degree four, moderate explicit
coefficients, and a proof-calculus interpretation without the native
machine tableau. It is not a variable-count improvement over those
thirty-seven or ten-variable equations.

The direct baseline is the newest incoming
\file{ConZFC_Shared_Certificates_2026-10-24.zip}, rather than the older
16,036-natural-variable direct compiler. The incoming archive contains
a completed $3,614$-natural-variable quartic, $2,954$ quadratic rows,
$15,996$ monomials and coefficient height $127,778,545$.
Its global inventory is $55$ power graphs, $16$ binomial
graphs and two factorial graphs \cite{shared}.

\begin{table}[htbp]
\centering\small
\begin{tabular}{lrrrr}
\toprule
Natural quartic & Variables & Rows & Monomials & Height\\
\midrule
Incoming direct & 3,614 & 2,954 & 15,996 & 127,778,545\\
Common radix & 2,108 & 1,745 & 10,629 & 127,777,275\\
Plus shared offset & 1,968 & 1,717 & 10,545 & 127,777,275\\
Shared Pell frame & 1,153 & 1,021 & 9,234 & 127,778,299\\
Compact natural & 667 & 535 & 23,747 & 2,649,643,818,040\\
Fewer monomials & 680 & 548 & 18,302 & 690,871,812,080\\
\bottomrule
\end{tabular}
\caption{Exact measurements for the direct natural-variable construction.
Every polynomial has degree four. The final integer companion has
\IntVars\ variables and \IntRows\ quadratic rows.}
\end{table}


Height means the maximum absolute coefficient of the expanded
polynomial. The support cap in the elimination pass exposes a tradeoff:
removing more coordinates can increase the number of monomials and
their coefficients. The low-height checkpoint is retained for that
reason. Every listed new point is a completed fixed polynomial;
no count describes only a fragment of an unfinished compiler.

The second improvement keeps the incoming thirteen-integer norm
architecture and its whole fixed coefficient tuple. Moving a square
multiplier inside one positivity certificate lowers its exact degree
from $1088$ to $1072$ at unchanged arithmetic cost.

The baseline implementation allocates its Pell environment separately
at every power call. The shared-power theorem gives a concrete
improvement over that inspected implementation. No worldwide
priority claim is made for sharing Diophantine witnesses.

\subsection{An explicit polynomial}

The primary natural equation is defined by the complete residual lists in
\file{direct/output/compact667/equation.json}. If those fixed lists are
$\mathcal L_0,\ldots,\mathcal L_{\NatRows-1}$, the equation is exactly
\begin{equation}
\boxed{\quad
 P_{\mathrm{batch}}(X_0,\ldots,X_{\NatVars-1})
 =
 \sum_{\ell=0}^{\NatRows-1}
 \left(\sum_{(c,\mu)\in\mathcal L_\ell}
           c\prod_{j\in\mu}X_j\right)^2=0.
 \quad}
\label{eq:literal}
\end{equation}
Each $c$ is a supplied integer and each $\mu$ is a sorted list of
at most two coordinate indices. The empty list is the constant
monomial; a repeated index denotes a square. The variable list
fixes the ordering.

This finite notation specifies an ordinary polynomial with fixed
integer coefficients. Its collected coefficient dictionary is also
included as \file{direct/output/compact667/expanded_polynomial.json}.
No variable exponent, factorial, binomial, remainder, bounded
quantifier or unnamed numerical parameter survives. The builder
reproduces data already delivered, rather than being the only
specification of an as-yet-uncomputed polynomial.

The integer equation has the same two representations in
\file{integer}. The thirteen-integer equation is a division-free
straight-line polynomial circuit in \file{guard13}; its fixed
leaves and inherited coefficient provenance are included.

\section{The compression principle}
\label{sec:principle}

An existential arithmetic certificate often chooses auxiliary numbers
only to make them sufficiently large, divisible or congruent. Compiling
each occurrence separately repeats those choices. We ask whether a
single choice can serve an entire finite computation.

Pointwise existence does not by itself justify this:
$\forall i\,\exists a\,R_i(a)$ does not imply
$\exists a\,\forall i\,R_i(a)$. We prove simultaneous completeness
for the particular certificates here. Every binomial digit can use
one sufficiently large radix. Every power graph can use one
dominating Pell parameter. A finite matrix period meets all of the
Pell-ordinate divisibility requirements at once.

The dominating parameter includes the proposed power outputs:
\[
 W=1+\sum_i(n_i+k_i+m_i).
\]
At an arbitrary polynomial zero, the global Pell norm forces
$A>W$, and the local modulus $T_i=2An_i-n_i^2-1$ then exceeds
$m_i$. Consequently the separate strict output-bound
certificates disappear.

There is no self-referential bound. The semantic arithmetic
values are chosen before the Pell witnesses. The new Pell
coordinates and the common congruence offset do not occur in
$W$. Numerically huge witnesses are compatible with a small
fixed number of quantified coordinates.

\section{One radix for every binomial coefficient}
\label{sec:radix}

For finitely many binomial requests $b_j=\binom{a_j}{r_j}$,
put $H_0=\sum_j a_j$ and
\[
 Q=2^{H_0}+1,\qquad L_j=Q^{r_j},\qquad C_j=(Q+1)^{a_j}.
\]
Every actual upper and lower argument is positive; in
particular $H_0>0$. Since $a_j\le H_0$, every coefficient
in the binomial expansion of $C_j$ is at most $2^{a_j}<Q$.
Its base-$Q$ digit in position $r_j$ is therefore exactly
$\binom{a_j}{r_j}$, including zero if $r_j>a_j$.

Four natural witnesses $h_j,u_j,v_j,w_j$ impose that digit:
\begin{equation}
 C_j=L_j(Qh_j+b_j)+u_j,\qquad
 L_j=u_j+v_j+1,\qquad Q=b_j+w_j+1.
\label{eq:shared-digit}
\end{equation}
The last two equations are the exact remainder and digit
bounds. Uniqueness of quotient and remainder proves the
equivalence.

The baseline has sixteen binomial graphs, each using three
powers. The shared construction uses sixteen numerator powers,
four distinct denominator powers, and one radix power.
Seven other powers remain unchanged. Thus the total inventory
falls from $55$ to $28$ power requests before their Pell
certificates are shared.

Some upper arguments depend on earlier factorial or binomial
outputs. The builder allocates those semantic outputs first,
records all requests, and only then emits the common radix.
For completeness choose the intended arithmetic values first;
those determine all upper arguments. For soundness an
arbitrary polynomial zero already assigns them values, and
the power theorem followed by the digit theorem forces their
identities. The radix is an auxiliary certificate parameter,
not part of the definition of the binomial function.

\section{One offset for signed congruence quotients}
\label{sec:offset}

The baseline represents $L_i\equiv R_i\pmod{D_i}$ by
$L_i+D_i a_i=R_i+D_i b_i$, with two natural quotient
coordinates. A fixed finite family needs only one common
natural offset $b$:
\begin{equation}
 L_i+D_i a_i=R_i+D_i b\qquad\text{for every }i.
\label{eq:shared-offset}
\end{equation}
The forward implication is immediate. Conversely write
$L_i-R_i=D_i t_i$ with signed integer $t_i$, choose
$b\ge\max(0,t_1,\ldots,t_s)$, and set $a_i=b-t_i$.
If $D_i=0$, the congruence means $L_i=R_i$ and places
no restriction on the offset. No coprimality hypothesis
is involved.

The common offset is chosen last and excluded from the
power-dominating sum. The construction imposes no
simultaneous congruence on the offset: it is simply an
upper bound on finitely many signed quotients.

\section{One Pell certificate for all positive powers}
\label{sec:batched-powers}

The principal new compression is at the level of existential certificates.
Every requested power retains its own input and output, but the auxiliary
number-theoretic environment is shared by the entire finite computation.
This is stronger than reusing identical arithmetic subexpressions: it
restricts independently chosen old witnesses to a common family, whose
simultaneous completeness must be proved.

For an integer $a\geq2$, let $\chi_a(t),\psi_a(t)\in\mathbb N$ be the Pell
coordinates defined by
\[
 \chi_a(t)+\psi_a(t)\sqrt{a^2-1}
       =(a+\sqrt{a^2-1})^t.
\]
These functions occur only in the proof. The emitted equation contains
ordinary polynomial expressions and no Pell-function symbol. Write
\[
 \mathcal N(a,x,y):\quad x^2+y^2=1+a^2y^2.
\]
All variables in this section range over $\mathbb N=\{0,1,2,\ldots\}$.

\subsection{The two established interfaces}

We use the positive-index case of the classical Pell-index representation.
For $a\geq2$ and $k>0$, the conjunction
\begin{align}
 &k\leq Y,\qquad \mathcal N(a,X,Y),\qquad
   \mathcal N(a,U,V),\qquad \mathcal N(B,S,R),\label{eq:index-norms}\\
 &B\geq2,\quad B\equiv1\pmod{4Y},\quad B\equiv a\pmod U,
   \quad V>0,\quad Y^2\mid V,\label{eq:index-frame}\\
 &S\equiv X\pmod U,\qquad R\equiv k\pmod{4Y}
 \label{eq:index-congruences}
\end{align}
is equivalent to $X=\chi_a(k)$ and $Y=\psi_a(k)$. In particular the
condition $k\leq Y$ makes $Y>0$ before any congruence modulus is used.

The positive-power representation says that, for $n,k>0$, $m=n^k$ is
equivalent to the existence of $W,a,T,Z,X,Y$ with $a\geq2$ satisfying
\begin{align}
 &X=\chi_a(k),\qquad Y=\psi_a(k),\label{eq:power-pell}\\
 &X+Yn\equiv Ya+m\pmod T,\qquad 2an=T+n^2+1,\label{eq:power-local}\\
 &m<T,\qquad n\leq W,\qquad k\leq W,
   \qquad \mathcal N(W+1,a,WZ).\label{eq:power-global}
\end{align}
These are the interfaces formalized as \texttt{Pell.matiyasevic} and
\texttt{Pell.eq\_pow\_of\_pell} in Mathlib's
\texttt{PellMatiyasevic.lean}~\cite{pell}. The shifted congruence in
\eqref{eq:power-local} avoids truncated subtraction. The standard proof
derives $a\geq n$, so it is equivalent to the source formulation
$X\equiv Y(a-n)+m\pmod T$.

\subsection{The batched polynomial conditions}

Consider a fixed nonempty list of power requests
\[
 (n_i,k_i,m_i),\qquad 1\leq i\leq s,
\]
whose bases and exponents are positive. Group identical exponent
expressions into $e_1,\ldots,e_g$, writing $k_i=e_{j(i)}$. Equality of
expressions is used here, not agreement at sampled numerical inputs.
Let
\[
 W=1+\sum_{i=1}^s(n_i+k_i+m_i).
 \label{eq:shared-W}
\]
Duplicate expressions may be included just once, provided the resulting
natural sum remains strictly larger than every base, exponent, and output.
Introduce a single frame $A,Z,U,V,B$, one quadruple
$(X_j,Y_j,S_j,R_j)$ for each exponent group, and one local modulus $T_i$
for each power request. Put
\[
 P=\prod_{j=1}^gY_j.
\]
The global conditions are
\begin{align}
 &A\geq2,\quad \mathcal N(W+1,A,WZ),\quad
   \mathcal N(A,U,V),\label{eq:batch-global-norms}\\
 &V=P^2(d+1),\qquad B=1+4Pq,
   \qquad B\equiv A\pmod U.\label{eq:batch-pooled}
\end{align}
Here $d,q$ are natural witnesses. For each exponent group retain
\begin{align}
 &e_j\leq Y_j,\quad \mathcal N(A,X_j,Y_j),\quad
     \mathcal N(B,S_j,R_j),\label{eq:batch-group-norms}\\
 &S_j\equiv X_j\pmod U,
       \qquad R_j\equiv e_j\pmod{4Y_j}.\label{eq:batch-group-congruences}
\end{align}
For each request retain only
\begin{equation}
 2An_i=T_i+n_i^2+1,\qquad
 X_{j(i)}+Y_{j(i)}n_i
      \equiv Y_{j(i)}A+m_i\pmod{T_i}.
 \label{eq:batch-local}
\end{equation}
The separate output-bound, $W$-bound, $B$-bound, local $V$-divisibility,
and local $B$-congruence certificates have disappeared.

\begin{theorem}[Simultaneous power certificate]
\label{thm:simultaneous-powers}
For every finite nonempty list of positive bases and exponents, the
conditions \eqref{eq:batch-global-norms}--\eqref{eq:batch-local} have
natural witnesses if and only if $m_i=n_i^{k_i}$ for every $i$.
If at least one base is known to be at least two, the explicit guard
$A\geq2$ may be omitted.
\end{theorem}

\begin{proof}[Soundness]
First suppose $A\geq2$. The last norm excludes $Z=0$, since that would
give $A=1$. Thus $Z\geq1$, and
\[
 A^2=1+(W^2+2W)W^2Z^2>W^2.
\]
Consequently $A>W$, which bounds every input and proposed output. Each
$Y_j\geq e_j>0$, so $P>0$ and
\[
 V=P^2(d+1)\geq Y_j^2\geq Y_j>0.
\]
In particular $Y_j^2\mid V$ and $B\equiv1\pmod{4Y_j}$.

The omitted guard $B\geq2$ follows from the retained equations. We have
$B=1+4Pq\geq1$. If $B=1$, then $\mathcal N(B,S_j,R_j)$ gives $S_j=1$.
The two norms with parameter $A$, together with $V\geq Y_j$, give
$U\geq X_j\geq2$. The congruence $S_j\equiv X_j\pmod U$ would then
imply $U\mid X_j-1$, although $0<X_j-1<U$. This is impossible.

All hypotheses of the Pell-index representation now hold group by group,
so $X_j=\chi_A(e_j)$ and $Y_j=\psi_A(e_j)$. For the omitted output bound,
$A\geq n_i+1$ yields
\[
 T_i-A=(2n_i-1)A-n_i^2-1
       \geq(n_i+2)(n_i-1)\geq0.
\]
Hence $m_i<W<A\leq T_i$. The standard positive-power representation
applies separately to every request and proves $m_i=n_i^{k_i}$.

Finally suppose the guard $A\geq2$ was omitted, but some $n_i\geq2$.
The retained modulus definition gives
$2An_i\geq n_i^2+1>2n_i$, and therefore $A>1$ before any preceding
argument is used. In the actual compiler a literal base-two request gives
the particularly simple equation $4A=T_i+5$.
\end{proof}

\begin{proof}[Simultaneous completeness]
Assume $m_i=n_i^{k_i}$ for every request, and first form the finite natural
number $W$ from these semantic inputs and outputs. Choose
\[
 A=\chi_{W+1}(W),\qquad
 Z=\frac{\psi_{W+1}(W)}{W}.
\]
The quotient is an integer because
$\psi_{W+1}(t)\equiv t\pmod W$. Also $W>0$, $A\geq2$, and
$\mathcal N(W+1,A,WZ)$ holds. For each $i$,
\[
 n_i^{k_i}\leq n_i^W<(W+1)^W\leq A.
\]
The standard Pell power congruence therefore supplies a positive
$T_i=2An_i-n_i^2-1>m_i$ and the local congruence
\eqref{eq:batch-local}. This is the source proof of the power
representation with a common upper bound in place of $\max(n_i,k_i)$.

Set $X_j=\chi_A(e_j)$ and $Y_j=\psi_A(e_j)$, so all $Y_j>0$, and form
$P=\prod_jY_j$. In the finite ring $\mathbb Z/(4P^2)$, the matrix
\[
 M=\begin{pmatrix}A&A^2-1\\1&A\end{pmatrix}
\]
is invertible because its determinant is one. Its powers belong to a
finite group. Some positive power is therefore the identity. Taking a
sufficiently large multiple $L$ of that period ensures $L>\max_j e_j$
and $L>1$. Put $U=\chi_A(L)$ and $V=\psi_A(L)$. The Pell matrix formula
gives
\[
 U\equiv1\pmod{4P^2},\qquad V\equiv0\pmod{4P^2}.
\]
Thus $V>0$, $P^2\mid V$, $U>A$, and $\gcd(U,4P)=1$.
The Chinese remainder theorem gives $B\geq2$ with
\[
 B\equiv A\pmod U,\qquad B\equiv1\pmod{4P}.
\]
Indeed the least nonnegative solution is already at least $A$, since
$A<U$. Define $d=V/P^2-1$ and $q=(B-1)/(4P)$.

For each group choose $S_j=\chi_B(e_j)$ and $R_j=\psi_B(e_j)$.
The recurrence with parameters congruent modulo $U$ yields
$S_j\equiv X_j\pmod U$. The recurrence modulo $B-1$ yields
$R_j\equiv e_j\pmod{4Y_j}$. All norms and bounds hold by construction,
so every retained condition is satisfied simultaneously.
\end{proof}

The finite-group step is deliberately stated for modulus $4P^2$.
Separately proving $Y_j^2\mid V$ for every $j$ would not prove
$P^2\mid V$, because the $Y_j$ can share prime factors. The common
period argument supplies the stronger divisibility actually used by the
emitted certificate.

\subsection{Compatibility with a common binomial radix}

The direct compiler also replaces the separate digit-extraction radices of
its finitely many binomial graphs by one radix. Given the positive upper arguments at the actual call sites,
$a_1,\ldots,a_r\in\mathbb N$, choose
\[
 H=\sum_{\ell=1}^{r}a_\ell,\qquad Q=2^H+1.
\]
Then $Q>2^{a_\ell}$ for every $\ell$. Every coefficient in
$(Q+1)^{a_\ell}$ is smaller than $Q$, since
$\binom{a_\ell}{j}\leq2^{a_\ell}$. Consequently quotient and remainder
equations for the digit at position $b_\ell$ recover
$\binom{a_\ell}{b_\ell}$ exactly, including the value zero when
$b_\ell>a_\ell$. Repeated lower arguments share the same power
$Q^{b_\ell}$. This is an ordinary positional-notation identity, with no
probabilistic or approximate digit extraction.

This common-radix construction and Theorem~\ref{thm:simultaneous-powers}
are compatible in a definite order. First choose a semantic solution of
the parent certificate and the common radix. Its finite power-request
list then has definite natural inputs and correct outputs. Next form
$W$ including all those outputs and choose the common Pell frame.
No Pell-frame coordinate enters the definition of the common radix or
the list of semantic power requests. Conversely, any zero of the final
equation yields all these natural values at once; the batched theorem
forces every requested power to have its intended value before the
binomial digit argument is invoked.

Both implications concern all natural numbers, including arbitrarily
large witnesses. A search over a finite numerical range is neither used
nor offered as a substitute for these proofs. The argument consists of
finite arithmetic, elementary Pell facts, finite-group periodicity, and
the Chinese remainder theorem, all formalizable in ZFC. Composing it
with the parent's direct C222 calibration preserves the stated
ZFC-provable consistency equivalence; it does not establish a new
kernel-checked or object-language ZFC derivation by itself.

\subsection{Eliminating terminal quadratic coordinates}

The final simplification acts on the literal quadratic residual system,
independently of the Pell implementation. Suppose one residual is
\[
 v-f(\mathbf x)=0,
\]
where $f\in\mathbb N[\mathbf x]$ has degree at most two and does not
contain $v$. If $f$ is affine, substitute it for $v$ everywhere. If $f$
is quadratic, make the same substitution provided $v$ occurs only
linearly in every other residual. Every new residual still has degree
at most two. Deleting the defining row and the variable $v$ gives an
exact bijection of natural zero sets: projection in one direction and
the unique restoration $v=f(\mathbf x)\in\mathbb N$ in the other.
The nonnegative coefficients are a sufficient, explicitly checked
condition for preserving the natural domain; a signed linear-algebra
elimination would not have this property automatically.

Repeated application provides a compact certificate consisting of the
deleted row, variable, and restoration polynomial at each step. Reading
the certificate backwards reconstructs every removed coordinate. The
verification checks exact polynomial identities, the domain condition,
and degree preservation at every step. It is stronger than comparing
the final polynomials at a selection of sample inputs.

\section{The direct proof-calculus source}
\label{sec:source}

This section records the inherited source mechanism
\cite{shared,flow} so that the consistency interpretation
is explicit. The new compression acts on the whole
source, not merely a local arithmetic relation.

The fixed numerical calculus $C_{222}$ uses implication
code $I(a,b)=2(a+b)(a+b+1)+4b$, with offsets $1,2,3$
for membership, universal quantification and equality.
Negation is $I(a,1)$. Appendix~\ref{app:calculus}
gives all constructors and twelve rule branches.
Truth is the freely available code $0$; contradiction
is code $1$. The numerical calculus is calibrated
to syntactic ZFC refutability by the named inherited
theorems in Section~\ref{sec:verification}.

\subsection{Finite flow certificates}

A padded rule with premises $p,q$ and conclusion $v$
has vector $e_v-e_p-e_q$, where $e_0=0$ and positive
labels are independent basis vectors. A finite list
of valid instances satisfying
\begin{equation}
 \sum_i(e_{v_i}-e_{p_i}-e_{q_i})=e_1
\label{eq:flow}
\end{equation}
already implies derivability of $1$.
Let $D$ be the closure under the listed rules of
truth and their axiom conclusions. If $1\notin D$,
sum the coordinates outside $D$. A conclusion
outside $D$ has at least one premise outside $D$,
so every rule contributes at most zero there.
The right side contributes one, a contradiction.
Only finitely many labels occur, so this is a
finite sum. Cycles and repeated premises cause
no exception.

Conversely a finite proof tree gives the flow by
counting its nodes. Expand it starting from the
open obligation $1$, replacing each conclusion
by its selected premises.

\subsection{A carrier recovering the exact flow}

Let $t,m\ge1$ be the number of instances and a
bound on their labels, and put $M=(m+1)^t$.
Set
\[
 A(X)=\prod_{i<t}(1+v_iX),\qquad
 C(X)=(1+X)\prod_{i<t}(1+p_iX)(1+q_iX).
\]
Both have nonnegative coefficients and constant
term one. Then
\begin{equation}
 A(M)=C(M)\quad\Longleftrightarrow\quad A(X)=C(X).
\label{eq:carrier}
\end{equation}
If $\deg C\ge t+1$, its evaluation is at least
$M^{t+1}$, whereas
\[
 A(M)\le(1+mM)^t<((m+1)M)^t=M^{t+1}.
\]
Thus each product has at most $t$ nonconstant
factors and coefficient sum at most $M$.
Every individual coefficient is strictly below
$M$, because the constant term is one and
a nonconstant product has another positive
coefficient. Every coefficient of $A-C$
therefore has absolute value below $M$.
Its least nonzero coefficient, if any, would
be divisible by $M$, by reducing its zero
evaluation after removing the corresponding
power of $M$. This is impossible.

Unique factorization in $\mathbb Q[X]$ now
recovers the exact multiplicities of the
distinct factors $1+jX$ for $j>0$, giving
\eqref{eq:flow}. No coprimality of their
integer values is assumed.

Natural accumulator states satisfy
\begin{equation}
 S_0=M+1,\quad S_t=1,\quad
 S_{i+1}(1+v_iM)=S_i(1+p_iM)(1+q_iM).
\label{eq:accumulator}
\end{equation}
Positive factors force every state positive,
and telescoping gives \eqref{eq:carrier}.
For completeness take $S_i$ to be the product
of the weights of current open obligations
in the proof-tree expansion. Every update
is then an integer identity.

One beta code represents the accumulator:
$\beta(x,y,i)=x\bmod(1+(i+1)y)$.
Choose $y$ sufficiently large and divisible
by every relevant positive index difference.
The moduli are pairwise coprime and exceed
all sequence entries, so CRT supplies $x$.
The choice order is finite proof, $t,m$,
carrier, accumulator, then $x,y$.
The later witness bound is independent of
the earlier label bound $m$.

\subsection{Thirteen local coordinates}

The local row uses three labels $p,q,v$,
four schema coordinates $a,b,c,d$ and six
beta/accumulator coordinates $S,S',u,r,u',r'$.
For $D_i=1+(i+1)y$ its arithmetic rows are
\begin{align*}
 x&=S+D_i u,&D_i&=S+r+1,\\
 x&=S'+(D_i+y)u',&D_i+y&=S'+r'+1,\\
 S'(1+vM)&=S(1+pM)(1+qM).
\end{align*}
The first four rows give exact consecutive
beta digits, since the divisors are positive.

For each rule branch square and sum its
residuals; multiply the twelve branch sums
to impose their disjunction. Add the
squares of the five other local rows.
This is a fixed nonnegative polynomial
$L(i,\mathbf z,\mathbf w)$.
Its implementation uses nonnegative-coefficient
circuits $E,F$ with $L=E-F$. For multiplication,
\[
 (E_1,F_1)(E_2,F_2)=
 (E_1E_2+F_1F_2,\ E_1F_2+F_1E_2).
\]
Addition and subtraction use componentwise
addition and the appropriate swap. The
difference is preserved and both components
are monotone in natural inputs.

\section{Removing the bounded universal quantifier}
\label{sec:crt}

Give the local witnesses bounds $u_j$ and
choose $V\ge u_j$ for every $j$. Define
\[
 H=E(t,\mathbf z,\mathbf u)+F(t,\mathbf z,\mathbf u)+t+V+1,
 \qquad T=H!,\qquad
 B_0=\prod_{i<t}(1+(i+1)T).
\]
Existence of bounded witnesses for every
index $i<t$ is equivalent to natural
$C,A_1,\ldots,A_d$ satisfying
\begin{align}
 B_0&\mid\binom{A_j}{u_j+1}
             &&\text{for each }j,\label{eq:crt-range}\\
 B_0&\mid1+(C+1)T,\label{eq:crt-index}\\
 E(C,\mathbf z,\mathbf A)&\equiv
       F(C,\mathbf z,\mathbf A)\pmod{B_0}.
       \label{eq:crt-row}
\end{align}

\begin{proof}[Soundness]
Fix $i<t$ and a prime divisor $p$ of
$D_i=1+(i+1)T>1$.
It is coprime to $T=H!$, so $p>H$.
The integer identity
\[
 (u_j+1)!\binom{A_j}{u_j+1}
       =\prod_{r=0}^{u_j}(A_j-r)
\]
also holds for $A_j\le u_j$, when both
sides vanish. Primality supplies a
$w_j\in[0,u_j]$ with $A_j\equiv w_j\pmod p$.
Comparing the index divisibility with
$D_i\equiv0$ and cancelling $T$ gives
$C\equiv i\pmod p$.
The row congruence yields
$E(i,\mathbf z,\mathbf w)\equiv
F(i,\mathbf z,\mathbf w)\pmod p$.
Monotonicity places both natural values
strictly below $p$, hence they are equal.
\end{proof}

\begin{proof}[Completeness]
A common divisor of $D_i,D_j$ for $i<j<t$
divides $(j-i)T$ and is coprime to $T$,
so it divides $j-i$. But $j-i\mid H!$,
making the common divisor one.
CRT now chooses $C\equiv i$ and
$A_j\equiv w_j^{(i)}\pmod{D_i}$
for valid bounded row witnesses.
Their falling-factorial products vanish
modulo every $D_i$. Since
$(u_j+1)!\mid T$ and $\gcd(B_0,T)=1$,
that factorial cancels, giving the
binomial divisibility. The index and
polynomial congruences follow by
substitution and pairwise coprimality.
\end{proof}

Soundness deliberately uses a prime
divisor of each modulus. A product
vanishing modulo a composite number
does not by itself supply a small root.
Completeness deliberately establishes
$H\ge u_j+1$ before factorial cancellation.

For the actual source, the three labels
have bound $m$ and the other ten local
coordinates have bound $V=m+U_0$ for an
extra natural $U_0$. There are thirteen
range-binomial graphs. Endpoint equations
directly encode $S_0=M+1,S_t=1$.
Writing $t=t_0+1,m=m_0+1$ loses no proof:
zero length contradicts the endpoints,
and zero label bound makes every update
constant. Encoded coordinates may be
shifted to $A_j=\alpha_j+1$, since adding
a multiple of $B_0$ preserves the
falling-factorial divisibility and its
factorial cancellation.

\subsection{Eliminating the remaining arithmetic macros}

The explanatory product is replaced by
\[
 Q_0=(1+(t+1)T)^t+2,\qquad
 B_0=
 \left(T^t\,t!\binom{Q_0+t}{t}\right)\bmod(TQ_0-1).
\]
The numerator is
$\prod_{i=1}^t(TQ_0+iT)$ and reduces
to the intended product modulo $TQ_0-1$.
That product is smaller than the
positive modulus, so it is the exact
remainder.

For a natural $n$, let
$\rho=(n+2)^{n+2}$,
$C=\binom{\rho+n}{n}$ and
$P=(\rho+1)^n$. Then
\[
 n!=\left\lceil P/C\right\rceil.
\]
For $n\ge2$,
\[
 \frac PC=n!\prod_{i=1}^n
 \left(1-\frac{i-1}{\rho+i}\right),
 \qquad
 0\le n!-\frac PC
 <\frac{n!\,n(n-1)}{2(\rho+1)}<1.
\]
The cases $0,1$ are immediate.
Natural $r,s$ impose the ceiling by
$P+C=Cf+r+1$ and $C=r+s+1$.
Each other division uses the analogous
quotient/remainder equations.
All binomials use the shared radix,
and all powers use the shared Pell
certificate. No arithmetic macro
remains in the final polynomial.

\subsection{The positivity order at arbitrary zeros}

All twenty-three exponent groups are
positive from structural shifts and
nonnegative arithmetic, before any
power output is known correct.
Consequently every group has
$Y_j\ge e_j>0$ and the common index
certificate recovers all its Pell
pairs first.

Twenty-seven of the twenty-eight
power bases are structurally positive.
The remaining base is the factorial
output $T$, which must be proved
positive rather than silently assumed.
Apply the local power argument first
to the structurally positive bases.
In particular the factorial numerator
$P=(\rho+1)^H$ is positive.
The retained ceiling rows
$P+C=CT+r+1$ and $C=r+s+1$
would imply $P\le0$ if $T=0$.
Hence $T>0$ and the final power
$T^t$ is valid.

The common radix power has literal
base two and a positive exponent, so
its correctness is available before
the binomial digit argument.
Its coefficient bound depends only
on the assigned natural upper
arguments, not on their prior
semantic correctness.
This order resolves the apparent
cycle introduced by deferring
factorial and binomial outputs.


\section{The literal quartic and its integer companion}
\label{sec:flatten}

Intermediate expressions use natural inputs, nonnegative
numeral leaves, addition and multiplication. Each signed
identity is stored as equality of two such expressions.
For a distinct arithmetic node introduce $z$ and impose
$z-x-y=0$ or $z-xy=0$. The gate graph is ordered, so its
values restore uniquely from semantic witnesses.
Squaring and summing its residuals produces one quartic.

The arithmetic gate graph remains acyclic when semantic
power outputs are allocated before their certificates:
those outputs are existential inputs of the graph,
not undefined gate values. Their intended meanings are
proved by the simultaneous theorem.

The inherited affine pass removes $z=L$ when $L$ is
affine with nonnegative coefficients. The new pass
also removes $z=F$ for a nonnegative quadratic $F$ if
$z$ occurs only linearly in every other remaining row.
Every deletion has a unique natural restoration and
preserves degree at most two in the residuals.
The complete certificate records every row, coordinate
and restoration polynomial. Its reverse order restores
all deleted coordinates.

\subsection{A safe unrestricted-integer translation}
\label{sec:integer}

For any natural coordinate $x$, the integer equation
\begin{equation}
 x=a^2+b^2+c(c+1)
\label{eq:three-square}
\end{equation}
is an exact nonnegativity guard. Its right side is
nonnegative for all integers $a,b,c$.
Conversely the three-square theorem represents $4x+1$;
modulo four exactly one of the three squares is odd
and the other two are even. Dividing by four gives
\eqref{eq:three-square}.

We guard the semantic coordinates of the shared-frame
checkpoint \emph{before} the final natural-coordinate
elimination. Its retained gate coordinates are then
nonnegative by topological induction from their
nonnegative arithmetic definitions. The integer
builder checks this entire dependency order.
Guarding only the coordinates named as witnesses after
a later elimination would require a new sign proof,
because a restored witness can depend on retained
gate values. We do not assume that inference.

After the guards have been added, monic substitutions
over $\Z$ can reduce the augmented quadratic system.
They preserve its exact integer zero set and all
remaining residual degrees. Their restoration
polynomials need not preserve positivity, since that
condition is already expressed by retained guard
equations and their substitutions.
The delivered integer polynomial has \IntVars\
coordinates, \IntRows\ quadratic rows and exact
degree four; it is supplied as residuals and as
collected coefficients.

\section{A secondary reduction: thirteen integers at degree 1072}
\label{sec:guard-supplement}

This section concerns the separate, large-coefficient Pell construction in
the incoming \file{ConZFC_13_Integer_Norms_2026-10-24.zip}.
It is independent of the direct compiler developed above.
The parent CRT construction has thirteen unrestricted integer variables,
exact degree $1088$, and $195$ live binary arithmetic operations.
One change to its joint positivity guard lowers the degree to $1072$,
with the same thirteen variables, the same $195$ operations, and exactly
the same six fixed integer coefficients. This is a degree-only improvement.

\subsection{Moving the nonzero scale inside the guard}

For integers $x,y,t$ put
\[
 \operatorname{Pair}(x,y,t)=xy-(t^2-x-y+1)^2,
 \qquad
 \operatorname{Pos}(p_1,p_2,p_3)
 =1+p_1^2+p_2^2+p_3(p_3+1).
\]
The two elementary interfaces are
\begin{align*}
 x>0\ \land\ y>0
 &\quad\Longleftrightarrow\quad
 \exists t\in\mathbb Z\ \operatorname{Pair}(x,y,t)>0,\\
 n>0
 &\quad\Longleftrightarrow\quad
 \exists p_1,p_2,p_3\in\mathbb Z\
 n=\operatorname{Pos}(p_1,p_2,p_3).
\end{align*}
For the first, positivity implies $xy>0$; two negative arguments make the
subtracted square exceed $xy$. For positive arguments, a nearest square
to $N=x+y-1$ has squared error strictly less than $N\le xy$.
The second follows from the three-square theorem applied to $4n-3$;
its residue modulo four gives one odd square and two even squares.

It follows that both equations
\begin{align}
 v^2\operatorname{Pair}(G,L,t)
   &=\operatorname{Pos}(p_1,p_2,p_3),\label{eq:guard-old}\\
 \operatorname{Pair}(v^2G,L,t)
   &=\operatorname{Pos}(p_1,p_2,p_3)\label{eq:guard-new}
\end{align}
express exactly $v\ne0$, $G>0$, and $L>0$ after existential
quantification of the four private guard variables. Their private
witnesses need not be the same. The benefit is algebraic: when $L$
has substantially higher degree than $v^2G$, the leading $L^2$ term
in~\eqref{eq:guard-new} avoids the outer factor $v^2$.

\subsection{Compatibility with the five-radical equation}

Retain the parent's thirteen coordinates
\[
 j,C,z,h,r,v,w,T,c,n,t,p_2,p_3\in\mathbb Z
\]
and every native, transport, Pell, and carrier definition. In particular,
\[
 J=6j+1,\quad q=\beta J+1,\quad W=12w+4,\quad
 X=qW,\quad Y=qv,\quad a=(X+1)Y,
\]
\[
 H=4a+3,\quad \Delta=a^2+H,\quad
 \eta=c-Yk,\quad \zeta=k-\eta,
\]
and
\[
 G=(A_{\rm cone}F)(B_{\rm cone}C)
    -(q-A_{\rm cone}F-B_{\rm cone}C)^2.
\]
Replace only the guard radicand by
\begin{equation}
 A_{\rm g}^{\rm new}
 =v^2G\eta\zeta-(t^2-v^2G-\eta\zeta+1)^2
    -1-p_2^2-p_3(p_3+1).
 \label{eq:guard-new-radicand}
\end{equation}
The single polynomial is the same cleared five-radical norm,
\begin{equation}
 \delta^{16}\mathcal F_{32}\!\left(
 Z;M^2\Lambda^2D_{\rm m},U^2D_f,U^2D_\tau,
 U^2A_{\rm g}^{\rm new},U^2C_y/\delta\right),
 \label{eq:guard-new-norm}
\end{equation}
where $M=2q^2$, $\Lambda=4H+1$, and $U=MH$.
Here $\mathcal F_{32}$ denotes the product over all thirty-two
independent signs. The supplied arithmetic graph clears the denominator
and eliminates the formal radicals; its operations are only $+,-,\times$.
The complete front end, coefficients, and literal graph accompany this
supplement, together with the complete unchanged parent package.

The parent's unconditional exceptional-case proof still gives
$U\ne0$, $\delta\ne0$, and four strictly positive non-guard radicands
at every integer assignment, including $v=0$ and $c=0$.
At a norm zero, a negative $A_{\rm g}^{\rm new}$ would therefore
contribute the only nonzero imaginary summand of a vanishing conjugate
factor. Thus $A_{\rm g}^{\rm new}\ge0$ and
\[
 \operatorname{Pair}(v^2G,\eta\zeta,t)
 =A_{\rm g}^{\rm new}+1+p_2^2+p_3(p_3+1)>0.
\]
The guard lemma recovers $v\ne0$, $G>0$, and $\eta\zeta>0$ in
the required order. In particular $c=0$ is impossible, because it would
give $\eta\zeta=-Y(Y+1)k^2\le0$.

The separation bound is unchanged. The parent proves $G<q^2$, while
$\eta\zeta\le k^2/4$, so
\[
 \sqrt{A_{\rm g}^{\rm new}}
 <|v|\sqrt{G\eta\zeta}<\frac{|v|\,|q|\,|k|}{2}.
\]
This is exactly its former bound for the guard radical. Consequently
the same strict ordering of radical magnitudes, rational-square
extraction, Pell descent, and two coprime congruences restore the same
calibrated source kernel.

For completeness start from a canonical source witness. Since $v^2G$
and $\eta\zeta$ are positive, choose a new $t$ by the guard lemma and
then a new positive representation with coordinates $p_1,p_2,p_3$.
This makes~\eqref{eq:guard-new-radicand} equal to $p_1^2$.
All other roots and source coordinates retain their values. The same
integral choice of center,
\[
 n=\Lambda\gamma+f+\tau+p_1+y-V-(v-q^2)/M,
\]
makes one norm factor vanish. Thus the new and old equations are
existence-equivalent over unrestricted integers. No identification of
their arbitrary private guard assignments is asserted.

\subsection{Exact degree, circuit size, and evidence}

The degree-seven leading part $\eta_0$ of $\eta$ gives the new
degree-twenty-eight guard component $-\eta_0^4$, replacing the old
degree-thirty component $-v^2\eta_0^4$. Hence the scaled guard
radicand $e=U^2A_{\rm g}^{\rm new}$ has degree forty rather than
forty-two. The independently expanded biquadratic cut has component
degree rows
\[
 (62,58,67,53),\qquad(134,120,129,125),\qquad
 (268,254,263,259).
\]
Its unique highest component is scalar. The final two quadratic norms
square it twice, giving exact total degree $4\cdot268=1072$.
An explicit coefficient on the actual fixed-coefficient slice is
\[
 [j^{480}w^{160}v^{192}h^{32}r^{64}c^{128}z^{16}]P
   =2^{976}3^{640}\beta^{480}>0.
\]
The exponent sum is $1072$, establishing attainment of the degree
bound independently of the other coefficient values.

The former multiplication $v^2\operatorname{Pair}(G,L,t)$ is
replaced by $v^2G$. The literal evaluator still has $41$ additions,
$122$ multiplications, and $32$ subtractions, totaling $195$.
Normal and optimized Python verification authenticate the whole
inherited source manifest, fixed numeral payloads, exact cut identities,
leading coefficient, and literal graph regeneration. Modular evaluations
against a separate conjugate-product evaluator are labeled transcription
diagnostics, rather than universal proofs.

The new theorem is this ordinary arithmetic existence reduction.
Its Con(ZFC) interpretation inherits the parent's named native/Pell and
K215 calibration arguments. The supplement supplies neither a new
end-to-end proof-assistant certificate nor a new object-language ZFC
derivation, and does not assert that the polynomial has no zero.


\section{Verification and the logical boundary}
\label{sec:verification}

The verification has three roles. Ordinary mathematics
proves the unbounded existence equivalences. Symbolic
audits bind the implemented source and arithmetic
graphs to those arguments. Exact coefficient
collection and certificate replay bind the final
polynomial to the graph. No finite check decides
whether the consistency equation itself has a zero.

The archive checks the fixed numerical calculus,
all shared-radix requests, every positive-power call
site, exact exponent grouping, the complete
dominating sum, the product of ordinates, the common
norms and congruences, every local modulus, and the
exclusion of certificate coordinates from the
dominating sum. It replays each elimination as an
integer-polynomial identity and independently expands
every delivered squared sum. Degree, coefficient
height, coordinate closure and variable use are
measured from the actual data.

Fixtures supplement those identities. They exercise
signed congruence quotients, binomial positions beyond
the degree, shared Pell periods and false power
outputs. They are small arithmetic examples, not
purported ZFC refutations. The receipts distinguish
normal and optimized Python runs. Required checks
raise explicit errors rather than depending on
assertions disabled by optimization.

\subsection{Inherited mathematical interfaces}

The direct source uses the calibrated equivalence
\[
 C_{222}(1)\longleftrightarrow
 \operatorname{Ref}_{\mathrm{ZF}}\longleftrightarrow
 \operatorname{Ref}_{\ZFC}.
\]
The named repository declarations are
\begin{quote}\small
\file{KEndpoint.c222_iff_zf_refutation}\\
\file{ConstructibleChoice.zf_refutation_iff_zfc_refutation}.
\end{quote}
They concern the fixed syntactic predicates. The
second uses constructible-universe equiconsistency;
neither assumes consistency. Relevant source copies,
notices and pins accompany the direct component.

The classical Pell-index and positive-power
representations are named interfaces in Mathlib
and the repository's literal representation modules
\cite{pell}. This contribution proves that their
witnesses can be selected simultaneously. It does
not reprove the full classical index theorem.

The thirteen-integer supplement has a separate
inherited native/Pell calibration. Packaging it
beside the direct quartic does not remove or replace
that dependency.

\subsection{Formalizability and reproducibility}

All new arguments use finite lists, polynomial
identities, elementary inequalities, Pell
recurrences, finite groups, finite CRT and finite
proof transformations. The matrix-period step is
a pigeonhole argument in a finite group. These
arguments are available in ZFC for its internal
natural numbers, with no restriction to standard
witnesses. Composing them with the identified
calibrations proves Theorem~\ref{thm:main}.

This is not a proof of $\Con(\ZFC)$ or a new
end-to-end formal-kernel certificate. Under the
usual consistency hypothesis, the no-zero
assertion is unprovable in ZFC by its equivalence
with consistency and the second incompleteness
theorem. No claim of two-sided independence from
mere consistency alone is needed.

The archive's \file{README.md} gives exact
reproduction commands. Its top-level
\file{verify.py} checks the delivered data and
rebuilds selected artifacts in temporary
directories. The TeX source compiles with
\texttt{pdflatex -interaction=nonstopmode
-halt-on-error report.tex}, run twice for
references. The contribution is prepared for
incorporation in ProveIt; it makes no remote
repository change.

\section{What this changes about the search}

The direct compiler's repeated cost was largely
the existential infrastructure for arithmetic
functions. Optimizing that infrastructure as a
finite family makes a substantial reduction
possible before any generic flattening.
The simultaneous-completeness proof is the
reusable result: sharing is justified by common
size bounds and finite periods, rather than
assumed from separate existential statements.

Further work could exploit proved relations
between different exponents, establish more
forced signs in the integer system, or translate
the finite certificates into a proof-assistant
development. None of those prospective gains
is included in the present counts.

The thirteen-integer equation retains very few
variables with large inherited constants. The
direct equation has degree four, moderate fixed
coefficients and a proof-calculus interpretation.
These are different useful improvements; no
comparison here proves a global simplest
Diophantine equation.

\appendix
\section{The fixed numerical calculus}
\label{app:calculus}

The following inherited numerical source
\cite{shared} makes the fixed logical
content of the coefficient dictionary
explicit. For natural inputs define
\begin{align*}
 I(a,b)&=2(a+b)(a+b+1)+4b,\\
 M(a,b)&=I(a,b)+1,& A(a,b)&=I(a,b)+2,& E(a,b)&=I(a,b)+3,\\
 N(a)&=I(a,1),&X(a,b)&=N(A(a,N(b))).
\end{align*}
They are numerical constructors, not
free function symbols in the final
polynomial. The eight unconditional
axiom-code families are
\begin{align*}
 r_1(a,b)&=X(a,E(a,b)),\\
 r_2(a,b,c)&=I(A(a,I(b,c)),I(A(a,b),A(a,c))),\\
 r_3(a,b,c)&=I(E(a,b),I(M(c,a),M(c,b))),\\
 r_4(a,b,c)&=I(A(a,A(b,c)),A(b,A(a,c))),\\
 r_5(a,b,c)&=I(E(c,a),I(M(c,b),M(a,b))),\\
 r_6(a,b)&=I(X(a,A(a,b)),b),\\
 r_G(a,b,c)&=I(E(c,a),I(E(c,b),A(M(a,b),E(a,b)))),\\
 r_8(a,b,c,d)&=I(I(I(c,d),a),I(I(a,c),I(b,c))).
\end{align*}
For the guarded bundle define
\begin{align*}
 R(b,k,h)&=A(b,I(h,N(A(b,I(h,N(M(b,k))))))),\\
 S(a,b,k)&=A(k,I(M(k,a),M(k,b))),\\
 C_0(a,b,k,h)&=I(S(a,b,k),R(b,k,h)),\\
 C_1(a,b,k,h)&=N(A(a,C_0(a,b,k,h))),\\
 C_2(a,b,k,h)&=N(A(a,I(h,C_1(a,b,k,h)))),\\
 C(a,b,k,h)&=A(b,I(M(b,k),C_2(a,b,k,h))),\\
 V(a,b,c)&=I(M(a,b),A(c,M(a,b))),\\
 K(a,b,c)&=A(c,I(C(a,b,c,c),N(R(b,c,c)))),\\
 \mathcal B(a,b,c)&=N(I(V(a,b,c),K(a,b,c))).
\end{align*}
Its side condition is $c\ne a,c\ne b$,
represented by one natural $d$ through
\[
 (c-a)^2(c-b)^2=(d+1)^2.
\]
For completeness take
$d=|(c-a)(c-b)|-1$.

The twelve padded branches for
premises $p,q$ and conclusion $v$ are
truth $p=q=v=0$; the eight axiom
branches $p=q=0,v=r_i$; the guarded
bundle $p=q=0,v=\mathcal B(a,b,c)$;
modus ponens $q=I(p,v)$; and
generalization $q=0,v=A(a,p)$.
The four schema coordinates are
shared across the disjunction,
but only the selected branch
must hold. Squaring each branch's
residuals and multiplying the
branch sums implements exactly
that convention.

The source audit matches the
composed constructors to the
pinned branch identities.
The calibration to ZFC
refutability uses the named
repository theorems; it is not
inferred merely from a resemblance
between these integers and
logical formulas.


\begin{thebibliography}{9}
\small\raggedright
\bibitem{repo}
Vladimir Reshetnikov and contributors,
\emph{ProveIt}, comparison snapshot \repoPin.
\url{https://github.com/VladimirReshetnikov/ProveIt/tree/a88aefd7ddab4771534bb049c67350e397c91b9a}.
\bibitem{latest}
ProveIt, \emph{Later arrivals and current size reductions},
\file{Computability/DiophantineConsistency/PROPOSAL_REVIEW_LATEST.md},
at the comparison snapshot.
\bibitem{wave11}
ProveIt, \emph{Wave11 review and size-reduction decisions},
\file{Computability/DiophantineConsistency/PROPOSAL_REVIEW_WAVE11.md},
at the comparison snapshot.
\bibitem{shared}
ProveIt incoming archive,
\file{ConZFC_Shared_Certificates_2026-10-24.zip},
\emph{Direct proof certificates and smaller Diophantine consistency equations}.
The archive filename is preserved as a source identifier.
Immediate source files, notices and hashes accompany this contribution.
\bibitem{flow}
ProveIt, \emph{A smaller complete source specification for direct proof flow},
\file{Computability/DiophantineConsistency/candidates/direct-flow-large-carrier/README.md},
at the comparison snapshot; and the inherited
\file{COMMUTATIVE_CERTIFICATES.md}.
\bibitem{pell}
Mario Carneiro and Mathlib contributors,
\emph{Pell's equation and Matiyasevic's theorem},
\file{Mathlib/NumberTheory/PellMatiyasevic.lean},
especially \file{Pell.matiyasevic} and
\file{Pell.eq_pow_of_pell}.
\url{https://leanprover-community.github.io/mathlib4_docs/Mathlib/NumberTheory/PellMatiyasevic.html}.
The archive includes the repository's pinned literal representation interfaces.
\bibitem{norm13}
ProveIt incoming archive,
\file{ConZFC_13_Integer_Norms_2026-10-24.zip}.
The unchanged parent package and its coefficient provenance
are included with the thirteen-integer supplement.
\end{thebibliography}
\end{document}
